\begin{frame}{Введение}
\begin{itemize}

    \item AlgoView — система трехмерной визуализации и интерактивного анализа информационных графов алгоритмов; AlgoLoad - платформа студенческих практикумов, использующая систему AlgoView.
    \item Цель работы - модернизация систем для улучшения функциональности и интерфейса, переработка устаревших подходов и методов обработки и визуализации графов.
    \item Технологическая платформа с AlgoView была улучшена и интегрирована в систему студенческих практикумов AlgoLoad;
    \item Система AlgoView была представлена на конференции ПаВТ'2024.
    
\end{itemize}
\end{frame}


\begin{frame}{Понятие информационного графа}

Информационный граф алгоритма — ациклический граф, вершины которого соответствуют операциям алгоритма, а дуги — связям по данным между этими операциями.  
 
\begin{itemize}
    % \item Информационный граф алгоритма — ациклический граф, вершины которого соответствуют операциям алгоритма, а дуги — связям по данным между этими операциями.
    \item Две вершины связываются дугой, если вторая использует данные, вычисленные в первой.
    \item Анализ информационного графа позволяет, например, определить множества независимых друг от друга операций, найти подходящее распределение операций по процессорам вычислительной системы, обнаружить узкие места.
\end{itemize}



% Информационный граф алгоритма — ациклический граф, вершины которого соответствуют операциям алгоритма, а дуги — связям по данным между этими операциями.

% Две вершины связываются дугой, если вторая использует данные, вычисленные в первой.

% Анализ информационного графа позволяет, например, определить множества независимых друг от друга операций, найти подходящее распределение операций по процессорам вычислительной системы, обнаружить узкие места.

    
\end{frame}

% длинное тире: —
